\section{Proofs for all-noise bounds, \texorpdfstring{$R_0^*$}{R0*}, and the lineage identity}\label{app:bounds}

\subsection{All-noise ring bounds, temporal independence, and host means}\label{subsec:app-bounds-any}

\begin{proof}[Proof of Proposition~{\upshape\ref{prop:bounds-any}}]
\begin{enumerate}
\item[(a)] Let $\boldsymbol{\Pi}_t = \mathbf{A}(t - 1)\cdots\mathbf{A}(0)$ and $\|\mathbf{X}\|_1 = \sum_{i,j} |X_{ij}|$. For each step $u$, $\|\mathbf{A}(u)\|_1 \le \sum_{k=1}^n (s_k + g_k h_k(u))$, while $\mathbf{1}^{\mathsf{T}}\mathbf{A}(u)\mathbf{x} \ge [\min_k g_k h_k(u)]\,\mathbf{1}^{\mathsf{T}}\mathbf{x}$ for every $\mathbf{x} \ge \mathbf{0}$. Because $\E|\ln h_k(0)| < \infty$, both $\ln^+\|\mathbf{A}(0)\|_1$ and $\ln^-[\min_k g_k h_k(0)]$ are integrable. Joint stationarity and ergodicity of $(h_1(t), \dots, h_n(t))$ therefore yield a finite deterministic limit $r = \lim_{t \to \infty} t^{-1}\ln\|\boldsymbol{\Pi}_t\|_1$ almost surely and in $L^1$ \citep{furstenberg1960}. Since $x_{\min}(0)\|\boldsymbol{\Pi}_t\|_1 \le X(t) = \mathbf{1}^{\mathsf{T}}\boldsymbol{\Pi}_t\mathbf{x}(0) \le x_{\max}(0)\|\boldsymbol{\Pi}_t\|_1$ whenever every $x_k(0) > 0$, $t^{-1}\ln X(t) \to r$ almost surely for every positive initial vector.

Choose any Markov chain on $\{1, \dots, n\}$ independent of the environment that stays in stage $k$ with probability $q_k$ ($q_k = 0$ if and only if $s_k = 0$, and $0 \le q_k < 1$ for all $k$) and moves to $k + 1$ with probability $1 - q_k > 0$. Let $\boldsymbol{\pi}$ be its stationary distribution and $Q(\omega)$ its stationary probability on $t$-step paths $\omega = (i_0, \dots, i_t)$ (for $n = 1$, treat stay and transfer as distinct labeled self-edges so $\omega$ records the edge sequence). Expanding the matrix product over paths gives
\[
X(t) = \sum_\omega Q(\omega)\,\frac{x_{i_0}(0)}{\pi_{i_0}} \prod_{u=0}^{t-1} \frac{a_{i_{u+1}i_u}(u)}{q_{i_u \to i_{u+1}}},
\]
where stay weights are $s_k$ and transfer weights are $g_k h_k(u)$. Applying Jensen's inequality to the concave logarithm for each environmental realization gives
\[
\ln X(t) \ge \sum_\omega Q(\omega)\ln\frac{x_{i_0}(0)}{\pi_{i_0}} + \sum_{u=0}^{t-1}\sum_{k=1}^n \pi_k \left[ q_k\ln\frac{s_k}{q_k} + (1 - q_k)\left(\ln\frac{g_k}{1 - q_k} + \ln h_k(u)\right) \right],
\]
omitting $q_k\ln(s_k/q_k)$ wherever $s_k = q_k = 0$. Divide by $t$ and let $t \to \infty$: the initial term vanishes, and Birkhoff's ergodic theorem gives $t^{-1}\sum_{u=0}^{t-1} \ln h_k(u) \to -\ell_k$ almost surely, so
\begin{equation}
r \ge \sum_{k=1}^n \pi_k \left[ q_k\ln\frac{s_k}{q_k} + (1 - q_k)\left(\ln\frac{g_k}{1 - q_k} - \ell_k\right) \right].
\label{eq:app-bounds-var}
\end{equation}
First set $q_k = s_k/\lambda$. Then $\tau_k = (1 - q_k)^{-1}$, $\pi_k = \tau_k/\mathcal{T}$, $\pi_k(1 - q_k) = 1/\mathcal{T}$, $\sum_{k=1}^n \pi_k q_k = (\mathcal{T} - n)/\mathcal{T}$, $\ln(s_k/q_k) = \ln\lambda$ whenever $s_k > 0$, and $\sum_{k=1}^n \ln[g_k/(1 - q_k)] = n\ln\lambda + \ln[\prod_k g_k / \prod_k (\lambda - s_k)] = n\ln\lambda$. Substituting into Eq.~\eqref{eq:app-bounds-var} yields $r \ge [(\mathcal{T} - n)\ln\lambda + n\ln\lambda - \sum_{k=1}^n \ell_k]/\mathcal{T} = \ln\lambda - \mathcal{T}^{-1}\sum_{k=1}^n \ell_k$. Next set $q_k = s_k$. Then $\pi_k = \tau_{k,0}/\mathcal{T}_0$, $\pi_k(1 - s_k) = 1/\mathcal{T}_0$, the staying log terms vanish, and Eq.~\eqref{eq:app-bounds-var} yields $r \ge \mathcal{T}_0^{-1}(\ln R_0 - \sum_{k=1}^n \ell_k)$. Dividing $m$ birth transfers by $R$ replaces $\ln R_0$ by $\ln R_0 - m\ln R$; evaluating at an existing root $R = R_0^*$ where $r = 0$ gives $0 \ge \ln R_0 - m\ln R_0^* - \sum_{k=1}^n \ell_k$, so $D \le \sum_{k=1}^n \ell_k$. When $\E h_k(0) = 1$, Jensen's inequality gives $-\ell_k = \E\ln h_k(0) \le \ln\E h_k(0) = 0$; for $h_k = \exp(\sigma_k z_k - \sigma_k^2/2)$ with $\E z_k = 0$, $\ell_k = \sigma_k^2/2$.
\item[(b)] For $n = 1$, $x_1(t) = x_1(0)\prod_{u=0}^{t-1} [s_1 + g_1 h_1(u)]$, so ergodicity gives $r = \E\ln(s_1 + g_1 h_1(0))$ directly, and dividing $g_1$ by $R_0^*$ with $R_0 = g_1/(1 - s_1)$ gives $g_1/R_0^* = (1 - s_1)\ee^D$. For any $n$ and any positive environment, $X(t) \ge x_k(t) \ge s_k^t x_k(0)$ for every $k$, giving $\liminf_{t \to \infty} t^{-1}\ln X(t) \ge \ln\max_k s_k$.
\item[(c)] When all $s_k = 0$, $\mathbf{A}(u)$ is a weighted cyclic permutation, so the $n$ initial stage cohorts never mix. Over $t = qn$ steps, the cohort starting in stage $j$ has log growth $\sum_{m=0}^{q-1} \sum_{r=0}^{n-1} [\ln g_{j+r} + \ln h_{j+r}(mn + r)]$. If the $n$-step shift is ergodic, $q^{-1}\sum_{m=0}^{q-1} \ln h_k(mn + r) \to -\ell_k$ almost surely for every $k$ and residue $r$, so every cohort (and hence their sum) grows at rate $n^{-1}\sum_{k=1}^n (\ln g_k - \ell_k) = \ln\lambda - n^{-1}\sum_{k=1}^n \ell_k$, giving equality in Eq.~\eqref{eq:bounds-growth} and $D = \sum_{k=1}^n \ell_k$. Under random-phase alternation between $(1.8, 0.2)$ and $(0.2, 1.8)$ at $n = 2, s_k = 0, g_k = 1$, the process is one-step ergodic (not $2$-step ergodic), and one cohort multiplies by $1.8$ at every step, giving $r = \ln 1.8 > 0 > \frac{1}{2}(\ln 1.8 + \ln 0.2)$. Finally, for $n = 2$ with $s_1 \in [0, 1)$, $s_2 = 0$, $h_1(t) = \exp(\sigma Z_t - \sigma^2/2)$, and $h_2(t) = \exp(-\sigma Z_{t-1} - \sigma^2/2)$, $x_2(t) = g_1 h_1(t - 1) x_1(t - 1)$ implies
\[
x_1(t + 1) = s_1 x_1(t) + g_2 h_2(t) x_2(t) = s_1 x_1(t) + g_1 g_2 \ee^{-\sigma^2} x_1(t - 1)
\]
because $h_2(t)h_1(t - 1) = \ee^{-\sigma^2}$ identically. This deterministic recurrence has growth rate zero if and only if $1 - s_1 = g_1 g_2 \ee^{-\sigma^2}$, so at threshold $g_1 g_2/(1 - s_1) = \ee^{\sigma^2}$ and $D = \sigma^2 = \ell_1 + \ell_2$ for every $s_1 \in [0, 1)$. \qedhere
\end{enumerate}
\end{proof}

\begin{proof}[Proof of Proposition~{\upshape\ref{prop:bounds-iid}}]
Since $\mathbf x(t)$ is $\mathcal F_t$-measurable, conditional expectation gives $\E\mathbf x(t+1)=\mathbf A_0\E\mathbf x(t)$ and hence $\E\mathbf x(t)=\mathbf A_0^t\E\mathbf x(0)$. Put $\lambda=\spr(\mathbf A_0)$. If $\lambda>0$, then for every $\epsilon>0$ there is $C_\epsilon<\infty$ such that $\mathbf1^{\mathsf T}\mathbf A_0^t\E\mathbf x(0)\leq C_\epsilon\exp\{t(\ln\lambda+\epsilon/2)\}$. Markov's inequality gives $\mathbb P(\|\mathbf x(t)\|_1>\exp\{t(\ln\lambda+\epsilon)\})\leq C_\epsilon\exp(-t\epsilon/2)$. Summing and applying Borel--Cantelli, then taking a countable sequence $\epsilon\downarrow0$, proves the asserted almost-sure upper bound. If $\lambda=0$, the finite matrix $\mathbf A_0$ is nilpotent, so $\mathbf x(t)=\mathbf0$ almost surely for all sufficiently large $t$. At a fixed negative birth tilt the additional hypothesis gives a strictly negative upper bound on growth, excluding a zero-growth root there. In the fixed-staying ring, reducing any designated positive birth transfer below its mean-threshold value makes the Perron root strictly less than one. Thus $\Delta\geq0$ and $D=m\Delta\geq0$ at any existing root.
\end{proof}

\begin{proof}[Proof of Remark~{\upshape\ref{rem:bounds-pairwise}}]
Set $\mathbf{A}_0 = \frac{1}{2}\mathbf{1}\mathbf{1}^{\mathsf{T}}$ ($\ln\lambda(\mathbf{A}_0) = 0$), $\mathbf{B} = \begin{psmallmatrix} 3/2 & -1/2 \\ -1/2 & -1/2 \end{psmallmatrix}$, and $\mathbf{A}(t) = \mathbf{A}_0 + \sigma\xi_t\mathbf{B}$ for $|\sigma| < 1/3$ (so $\mathbf{A}(t) > \mathbf{0}$). Concatenate independent three-step blocks $(X_j, Y_j, X_j Y_j)$ where $X_j, Y_j \in \{-1, +1\}$ are independent fair signs, with a uniformly random initial block phase in $\{0, 1, 2\}$. This driver $\xi_t$ is stationary and ergodic with $\E\xi_t = 0$, $\E[\xi_t\xi_s] = 0$ for all $t \ne s$, and $\E[\xi_t\xi_{t-1}\xi_{t-2}] = 1/3$. Writing the state in the basis $\mathbf{v} = (1/2, 1/2)^{\mathsf{T}}$, $\mathbf{u} = (1/2, -1/2)^{\mathsf{T}}$ as $\mathbf{y}_t = \mathbf{v} + z_t\mathbf{u}$ ($\mathbf{1}^{\mathsf{T}}\mathbf{y}_t = 1$), $\mathbf{A}_0\mathbf{v} = \mathbf{v}$, $\mathbf{A}_0\mathbf{u} = \mathbf{0}$, $\mathbf{B}\mathbf{v} = \mathbf{u}$, and $\mathbf{B}\mathbf{u} = \mathbf{v} + \mathbf{u}$ give
\[
G_t = \mathbf{1}^{\mathsf{T}}\mathbf{A}(t)\mathbf{y}_t = 1 + \sigma\xi_t z_t, \qquad z_{t+1} = \frac{\sigma\xi_t(1 + z_t)}{1 + \sigma\xi_t z_t}.
\]
Expanding in $\sigma$ gives $z_{t+1} = \sigma\xi_t + \sigma^2\xi_t\xi_{t-1}(1 - \sigma\xi_t) + O(\sigma^3)$ and $\ln G_t = \sigma\xi_t z_t - \frac{1}{2}\sigma^2 z_t^2 + O(\sigma^4)$. Since $\E[\xi_t\xi_{t-1}]=0$, the quadratic coefficient of $\E\ln G_t$ vanishes; the term $-\sigma^2\E z_t^2/2$ is of fourth order: $\E[\sigma\xi_t(\sigma\xi_{t-1})] = 0$, while at order $\sigma^3$, $\sigma\xi_t(\sigma^2\xi_{t-1}\xi_{t-2})$ has expectation $\sigma^3\E[\xi_t\xi_{t-1}\xi_{t-2}] = \sigma^3/3$, yielding $r(\sigma) = \sigma^3/3 + O(\sigma^4) > 0$ for small $\sigma > 0$.
\end{proof}

\subsection{Further results and proofs for Section~\ref{sec:bounds}}

\begin{corollary}[Arithmetic and geometric host means]\label{cor:bounds-means}
Under the fixed-staying ring hypotheses of Proposition~\ref{prop:bounds-any}, write each transfer as $g_k N_k(t)$ with $N_k(t) > 0$ jointly stationary and ergodic and $\E N_k(0), \E|\ln N_k(0)| < \infty$. Let $A_k = \E N_k(0)$ and $G_k = \exp(\E\ln N_k(0))$ be the arithmetic and geometric means of host $k$, and let $R_0^{\mathrm{A}}$ and $R_0^{\mathrm{G}}$ be the per-circuit reproduction numbers when hosts are fixed at $A_k$ and $G_k$. At any zero-growth divisor $R_0^*$ on $m$ counted births,
\begin{equation}
R_0^{\mathrm{G}} \le (R_0^*)^m, \qquad \text{and} \qquad (R_0^*)^m \le R_0^{\mathrm{A}} \quad \text{when } \mathbf{A}(t) \text{ are independent across time}.
\label{eq:bounds-means}
\end{equation}
\end{corollary}

\begin{proof}[Proof of Corollary~{\upshape\ref{cor:bounds-means}}]
Put $h_k(t) = N_k(t)/A_k > 0$, so $\E h_k(0) = 1$, the baseline transfers at arithmetic host means are $\bar{g}_k = g_k A_k$, and $\ell_k = -\E\ln[N_k(0)/A_k] = \ln A_k - \ln G_k$. Then $\ln R_0^{\mathrm{A}} - \sum_{k=1}^n \ell_k = \sum_{k=1}^n \ln[g_k G_k/(1 - s_k)] = \ln R_0^{\mathrm{G}}$. Proposition~\ref{prop:bounds-any}(a) gives $\ln R_0^{\mathrm{A}} - m\ln R_0^* = D \le \sum_{k=1}^n \ell_k$, i.e., $m\ln R_0^* \ge \ln R_0^{\mathrm{G}}$, and Proposition~\ref{prop:bounds-iid} gives $D \ge 0$, i.e., $m\ln R_0^* \le \ln R_0^{\mathrm{A}}$, when the matrices $\mathbf A(t)$ are independent across time.
\end{proof}

\subsubsection{Random-environment \texorpdfstring{$R_0^*$}{R0*}, path representation, and lineage variance}\label{subsec:app-bounds-lineage}

\begin{proof}[Proof of Proposition~{\upshape\ref{prop:bounds-r0star}}]
\begin{enumerate}
\item[(a)] Dividing the $m$ counted birth transfers by $R > 0$ changes the deterministic characteristic equation to $\prod_{k=1}^n (\lambda - s_k) = R^{-m}\prod_{k=1}^n g_k$, which has root $\lambda = 1$ if and only if $R^m = \prod_{k=1}^n g_k/(1 - s_k) = R_0$, proving $\spr(\mathbf{F}(\mathbf{I} - \mathbf{V})^{-1}) = R_0^{1/m}$. Writing $\Phi(\lambda) = \sum_{k=1}^n \ln[(\lambda - s_k)/(1 - s_k)] = \ln R_0$ at fixed stays $0 \le s_k < 1$, we have $\Phi(1) = 0$ and $\dd\Phi/\dd(\ln\lambda)|_{\lambda=1} = \sum_{k=1}^n (1 - s_k)^{-1} = \mathcal{T}_0 > 0$, so the inverse function theorem gives $\ln\lambda = \ln R_0 / \mathcal{T}_0 + O((\ln R_0)^2)$ as $\ln R_0 \to 0$ (and $\Phi(\lambda) = \ln\lambda$ identically when $n = 1, s_1 = 0$).
\item[(b)] For any constants $d_k > 0$, the diagonal change of units $\tilde{x}_k = d_k x_k$ leaves $s_k$ and $h_k(t)$ unchanged, replaces $g_k$ by $g_k d_{k+1}/d_k$, and changes $\|\mathbf{x}(t)\|_1$ by at most $\max_k d_k / \min_k d_k$, leaving $r$ unchanged. Since any two positive transfer vectors with the same product $P = \prod_{k=1}^n g_k$ are related by such a diagonal scaling, $r = r(y)$ depends on $g_1, \dots, g_n$ solely through $y = \ln P$. Dividing $m$ birth transfers by $R$ sends $y \mapsto y - m\ln R$. If $y_0$ is the unique root of $r(y) = 0$, then $r(\ln P - m\ln R_0^*) = 0$ gives $m\ln R_0^* = \ln P - y_0$ and $D = \ln R_0 - m\ln R_0^* = y_0 - \sum_{k=1}^n \ln(1 - s_k)$, which depends on neither $g_k$ nor $m$.

To prove existence and uniqueness of $y_0$, when $n = 1$ and $s_1 = 0$, $r(y) = y - \sigma_1^2/2$ has the unique root $y_0 = \sigma_1^2/2$. Otherwise at least one $s_k > 0$. Set $g_k = c_k \ee^{y/n}$ with $\prod_{k=1}^n c_k = 1$. Every $t$-step path with $j$ transfers contributes a positive weight times $\ee^{jy/n}$; a sum of positive log-linear exponentials is log-convex and nondecreasing in $y$, so $t^{-1}\ln X(t)$ and its limit $r(y)$ are finite, convex, and nondecreasing on $\mathbb{R}$, hence continuous. A path that stays forever in a stage with $s_k > 0$ gives $r(y) \ge \max_k \ln s_k$, while the column-sum bound $r(y) \le \E\ln\max_k(s_k + c_k\ee^{y/n}h_k(0)) \to \max_k \ln s_k < 0$ as $y \to -\infty$ by dominated convergence. A path that transfers at every step gives $r(y) \ge y/n + n^{-1}\sum_{k=1}^n (\ln c_k - \ell_k) \to +\infty$ as $y \to +\infty$. By continuity $r(y)$ has a zero $y_0$. If $r(y_1) = r(y_2) = 0$ for $y_1 < y_2$, monotonicity and convexity would force $r(y) = 0$ on $[y_1, y_2]$ and $r(y) \ge 0$ for all $y < y_1$, contradicting $\lim_{y \to -\infty} r(y) = \max_k \ln s_k < 0$.
\item[(c)] Shift $y$ to $y = \ln[\prod_{k=1}^n g_k / \prod_{k=1}^n (1 - s_k)]$, so deterministic threshold is $y = 0$ and $D(\sigma) = y_0(\sigma)$ solves $r(y_0(\sigma), \sigma) = 0$. At $y = 0$, Theorem~\ref{thm:ring-rem-ar1} gives $r(0, \sigma) = -C_0\sigma^2 + O(\sigma^4)$. For any fixed $\epsilon > 0$, $r(\pm\epsilon, \sigma) \to L(\pm\epsilon) = \ln\lambda(\pm\epsilon)$ as $\sigma \to 0$ with $L(-\epsilon) < 0 < L(\epsilon)$, so $y_0(\sigma) \to 0$. By convexity of $y \mapsto r(y, \sigma)$, for $0 < |y_0(\sigma)| < \epsilon$ the chord slope $[r(y_0(\sigma), \sigma) - r(0, \sigma)]/y_0(\sigma)$ lies between $[r(0, \sigma) - r(-\epsilon, \sigma)]/\epsilon$ and $[r(\epsilon, \sigma) - r(0, \sigma)]/\epsilon$. Letting $\sigma \to 0$ and then $\epsilon \downarrow 0$ sends both bounds to $L'(0) = 1/\mathcal{T}_0$, giving $C_0\sigma^2 + O(\sigma^4) = y_0(\sigma)[1/\mathcal{T}_0 + o(1)]$ and $D = \mathcal{T}_0 C_0\sigma^2 + o(\sigma^2) = \Gamma_0 + o(\sigma^2)$.

For the fourth-order bound, choose $\delta > 0$ so that $L'(y) \ge \ell > 0$ on $|y| \le \delta$. On $|y| \le \delta$, the resolvent expression for $C_y$ in Eq.~\eqref{eq:ring-rem-ar1} is smooth, so $|C_y| \le \bar{c}$ and $|C_y - C_0| \le M_1|y|$. Lemma~\ref{lem:ring-rem-uniform} gives $r(y, \sigma) = L(y) - C_y\sigma^2 + \rho(y, \sigma)$ with $|\rho(y, \sigma)| \le K\sigma^4$ uniformly on $|y| \le \delta, |\sigma| \le \sigma_0$. Shrinking $\sigma_0$ if needed so $r(-\delta, \sigma) < 0 < r(\delta, \sigma)$, the unique root $y_0$ lies in $(-\delta, \delta)$ and satisfies $L(y_0) = C_{y_0}\sigma^2 - \rho(y_0, \sigma)$. Thus $|L(y_0)| \le \bar{c}\sigma^2 + K\sigma^4$ and $|y_0| \le (\bar{c}\sigma^2 + K\sigma^4)/\ell = O(\sigma^2)$, whence $|C_{y_0} - C_0| \le M_1|y_0| = O(\sigma^2)$ and $L(y_0) = C_0\sigma^2 + O(\sigma^4)$. Since $L(y_0) = y_0/\mathcal{T}_0 + O(y_0^2) = y_0/\mathcal{T}_0 + O(\sigma^4)$, we conclude $D = y_0 = \mathcal{T}_0 C_0\sigma^2 + O(\sigma^4) = \Gamma_0 + O(\sigma^4)$. \qedhere
\end{enumerate}
\end{proof}

For general life cycles, resolve $\mathbf{A}_0$ into finitely many positive labeled edges $e: j \to i$ of mean weights $a_e > 0$. Keep parallel stays, transfers, routes, and movements distinct when their multipliers differ. The environment-independent reference lineage takes edge $e$ with probability $P(e: j \to i) = w_i a_e / (\lambda w_j)$, has stationary law $\pi_j = v_j w_j$, and uses that edge with frequency $c_e = w_i a_e v_j / \lambda$.

\begin{proposition}[Path representation and general visit-weighted bound]\label{prop:bounds-path}
Let $\mathbf{A}_0$ have eigenvalue $\lambda > 0$ and strictly positive left eigenvector $\mathbf{w}$ satisfying $\mathbf{w}^{\mathsf{T}}\mathbf{A}_0 = \lambda\mathbf{w}^{\mathsf{T}}$. For any nonnegative finite multipliers $h_e(t)$ and every integer $T \ge 1$,
\begin{equation}
[\mathbf{A}(T)\cdots\mathbf{A}(1)]_{ij} = \lambda^T \frac{w_j}{w_i}\,\E_{P,j}\!\left[\prod_{t=1}^T h_{e_t}(t)\,\mathbf{1}_{\{X_T = i\}}\right].
\label{eq:bounds-L1}
\end{equation}
Suppose further that $\mathbf{A}_0$ is irreducible (with $\mathbf{w}^{\mathsf{T}}\mathbf{v} = 1$), the multipliers $h_e(t) > 0$ are jointly stationary and ergodic with $\E\max_e|\ln h_e(0)| < \infty$, and the population starts with positive support in every state. Put $S_T = \sum_{t=1}^T \ln h_{e_t}(t)$ and $\ell_e = -\E\ln h_e(0)$. Then
\begin{equation}
r - \ln\lambda = \lim_{T \to \infty} \frac{1}{T}\,\E_{\mathrm{env}}\!\left[\ln\E_{P,\pi}\ee^{S_T}\right] \ge -\sum_e c_e \ell_e.
\label{eq:bounds-L2}
\end{equation}
When $\E h_e(0) = 1$, $\ell_e \ge 0$; if $\mathbf{A}_0$ is scaled to $\lambda = 1$, counted birth edges have total reference frequency $b > 0$, and multiplying those birth edges by $\ee^y$ has a finite zero-growth root $y = \Delta$, then $\Delta \le b^{-1}\sum_e c_e \ell_e$ (with $D = m\Delta$ in a fixed-staying ring of $m$ birth transfers). In a ring at $\lambda = 1$ with one birth transfer whose staying rates $S_k(t)$ also carry positive mean-one multipliers, $D \le \sum_{k=1}^n \ell_k^g + \sum_{k=1}^n (\tau_k - 1)\ell_k^s$ where $\tau_k = (1 - \E S_k)^{-1}$ and zero-mean stays are omitted.
\end{proposition}

\begin{proof}[Proof of Proposition~{\upshape\ref{prop:bounds-path}}]
Resolve $\mathbf{A}_0$ into finitely many positive labeled edges $e: j \to i$ of weights $a_e > 0$. Since $\sum_{e \text{ from } j} w_i a_e / (\lambda w_j) = (\mathbf{w}^{\mathsf{T}}\mathbf{A}_0)_j / (\lambda w_j) = 1$, $P(e: j \to i) = w_i a_e / (\lambda w_j)$ is a Markov transition kernel. Expanding $[\mathbf{A}(T)\cdots\mathbf{A}(1)]_{ij}$ over $T$-step edge paths $(e_1, \dots, e_T)$ from $j$ to $i$, the product of reference probabilities along a path is $\lambda^{-T}(w_i / w_j)\prod_{t=1}^T a_{e_t}$ because intermediate $w$-factors telescope, yielding Eq.~\eqref{eq:bounds-L1} immediately.

When $\mathbf{A}_0$ is irreducible, $\mathbf{v}, \mathbf{w} > \mathbf{0}$ with $\mathbf{w}^{\mathsf{T}}\mathbf{v} = 1$, $\pi_j = v_j w_j$ is stationary for $P$, and $c_e = \pi_j P(e: j \to i) = w_i a_e v_j / \lambda$. Summing Eq.~\eqref{eq:bounds-L1} over $i, j$ against $w_i$ and $v_j$ gives $\E_{P,\pi}\ee^{S_T} = \lambda^{-T}\mathbf{w}^{\mathsf{T}}\mathbf{A}(T)\cdots\mathbf{A}(1)\mathbf{v}$. Because $\mathbf{v}, \mathbf{w} > \mathbf{0}$, $\mathbf{w}^{\mathsf{T}}\boldsymbol{\Pi}_T\mathbf{v}$ is comparable up to fixed positive constants to $\|\boldsymbol{\Pi}_T\|_1$ and to $\mathbf{1}^{\mathsf{T}}\boldsymbol{\Pi}_T\mathbf{x}(0)$ for any positive-support initial state. Under joint stationarity, ergodicity, and $\E\max_e |\ln h_e(0)| < \infty$, Kingman's subadditive ergodic theorem gives $T^{-1}\ln\E_{P,\pi}\ee^{S_T} \to r - \ln\lambda$ almost surely and in $L^1$, proving the limit in Eq.~\eqref{eq:bounds-L2}. For each environment, Jensen's inequality gives $\ln\E_{P,\pi}\ee^{S_T} \ge \E_{P,\pi} S_T$. Because $(P, \pi)$ is stationary and independent of the environment, $\E_{\mathrm{env}}\E_{P,\pi} S_T = \sum_{t=1}^T \sum_e c_e \E\ln h_e(t) = -T\sum_e c_e \ell_e$; dividing by $T$ gives $r - \ln\lambda \ge -\sum_e c_e \ell_e$. At $\lambda = 1$, multiplying counted birth edges (total frequency $b > 0$) by $\ee^y$ adds $y$ to $\ln h_e$ on each birth edge, giving $r(y) \ge b y - \sum_e c_e \ell_e$; substituting an existing zero-growth root $y = \Delta$ yields $\Delta \le b^{-1}\sum_e c_e \ell_e$. In a mean-threshold ring with one birth transfer ($b = 1/\mathcal{T}_0$), transfer edges have $c_k^g = 1/\mathcal{T}_0$ and positive-mean staying edges have $c_k^s = \pi_k \E S_k = (\tau_{k,0} - 1)/\mathcal{T}_0$, giving $D = \Delta \le \sum_{k=1}^n \ell_k^g + \sum_{k=1}^n (\tau_{k,0} - 1)\ell_k^s$.
\end{proof}

\begin{theorem}[Second-order lineage-variance identity]\label{thm:bounds-lineage}
Let $\mathbf{A}_0$ be finite and primitive. Consider a smooth mean-preserving multiplier family satisfying $\ln h_e(\sigma, t) = \sigma f_e(t) + \sigma^2 g_e(t) + O(\sigma^3)$, where $(f_e(t), g_e(t))$ are jointly stationary and ergodic with finite second moments, $\E f_e(0) = 0$, $\E g_e(0) = -\frac{1}{2}\E[f_e(0)^2]$, and bounded lag covariances. Along a stationary reference line $(P, \pi)$, define $Y_T = \sum_{t=1}^T f_{e_t}(t)$, $J_2 = \frac{1}{2}\sum_e c_e \E[f_e(0)^2]$, and
\begin{equation}
V_2 = \lim_{T \to \infty} \frac{1}{T}\,\E_{\mathrm{env}}\!\left[\Var_{P,\pi}(Y_T \mid \text{environment})\right],
\label{eq:bounds-V2}
\end{equation}
where the limit is that of the finite-time total-variance difference. Write $J = \sigma^2 J_2$ and $V_{\mathrm{between}} = \sigma^2 V_2$ (so $J = \frac{1}{2}\sum_e c_e \sigma_e^2$ when $f_e = \kappa_e\xi_{\delta(e)}$ with unit-variance drivers and $\sigma_e = \kappa_e\sigma$). Then $J_2 - \frac{1}{2}V_2$ equals the lag-sum coefficient $C''$ of Eq.~\eqref{eq:ring-rem-lin}, and under uniform-in-time remainder control,
\begin{equation}
\ln\lambda - r(\sigma) = \sigma^2\left(J_2 - \frac{1}{2}V_2\right) + \text{remainder} = J - \frac{1}{2}V_{\mathrm{between}} + \text{remainder},
\label{eq:bounds-lineage}
\end{equation}
with remainder $O(\sigma^4)$ for mean-one lognormal multipliers driven by Gaussian linear processes satisfying Theorem~\ref{thm:ring-rem-lin}, and remainder $O(|\sigma|^3)$ for uniformly bounded centered additive perturbations $\mathbf{A}_0 + \sigma\xi(t)\mathbf{D}$ supported on positive entries of $\mathbf{A}_0$ (improving to $O(\sigma^4)$ when the full joint process $\xi$ is invariant in law under $\xi \mapsto -\xi$).
\end{theorem}

\begin{proof}[Proof of Theorem~{\upshape\ref{thm:bounds-lineage}}]
Write $\gamma_{ee'}(\ell) = \E[f_e(t)f_{e'}(t - \ell)]$ for $\ell \ge 0$, which is bounded by Cauchy--Schwarz and stationarity. For a fixed environment and finite $T$, differentiating $\Psi_T(\sigma) = \ln\E_{P,\pi}\ee^{S_T(\sigma)}$ ($S_T(\sigma) = \sum_{t=1}^T \ln h_{e_t}(\sigma, t)$) twice at $\sigma = 0$ gives $\Psi_T'(0) = \E_{P,\pi} Y_T$ and $\Psi_T''(0) = \E_{P,\pi}[\sum_{t=1}^T 2g_{e_t}(t)] + \Var_{P,\pi}(Y_T \mid \text{env})$. Taking expectations over the environment with $\E f_e(0) = 0$ and $\E g_e(0) = -\frac{1}{2}\gamma_{ee}(0)$ gives $\E_{\mathrm{env}}\Psi_T'(0) = 0$ and
\[
-\frac{1}{2T}\E_{\mathrm{env}}\Psi_T''(0) = J_2 - \frac{1}{2T}\E_{\mathrm{env}}\Var_{P,\pi}(Y_T \mid \text{env}).
\]

To evaluate $V_2$ without assuming separate convergence of the unconditional and environment-mean variance rates, apply the law of total variance at finite $T$, writing $\E_{\mathrm{env}}\Var_{P,\pi}(Y_T \mid \text{env})$ as $\Var_{\mathrm{total}}(Y_T)$ minus $\Var_{\mathrm{env}}(\E_{P,\pi}[Y_T \mid \text{env}])$. Along a stationary reference path, one edge is traversed at each step $t$, so $\E_{\mathrm{total}}[f_{e_t}(t)^2] = \sum_e c_e \gamma_{ee}(0)$, while at lag $\ell = t - t' \ge 1$ the cross-expectation is $\sum_{e,e'} c_{e'} u_{e' \to e}(\ell)\gamma_{ee'}(\ell)$. Meanwhile $\E_{P,\pi}[Y_T \mid \text{env}] = \sum_{t=1}^T \sum_e c_e f_e(t)$ has variance $T\sum_{e,e'} c_e c_{e'}\gamma_{ee'}(0) + 2\sum_{\ell=1}^{T-1} (T - \ell)\sum_{e,e'} c_e c_{e'}\gamma_{ee'}(\ell)$. Subtracting and dividing by $T$ gives
\[
\begin{aligned}
\frac{1}{T}\E_{\mathrm{env}}\Var_{P,\pi}(Y_T \mid \text{env}) &= \sum_e c_e\gamma_{ee}(0) - \sum_{e,e'} c_e c_{e'}\gamma_{ee'}(0) \\
&\quad + 2\sum_{\ell=1}^{T-1}\left(1 - \frac{\ell}{T}\right)\sum_{e,e'} \gamma_{ee'}(\ell) c_{e'}\bigl[u_{e' \to e}(\ell) - c_e\bigr].
\end{aligned}
\]
Because $\mathbf{A}_0$ is primitive, $|u_{e' \to e}(\ell) - c_e|$ decays geometrically in $\ell$; boundedness of $\gamma_{ee'}(\ell)$ therefore makes the infinite series converge absolutely, and dominated convergence as $T \to \infty$ gives $V_2 = \sum_e c_e\gamma_{ee}(0) - \sum_{e,e'} c_e c_{e'}\gamma_{ee'}(0) + 2\sum_{\ell=1}^{\infty} \sum_{e,e'} \gamma_{ee'}(\ell) c_{e'}[u_{e' \to e}(\ell) - c_e]$. Subtracting $\frac{1}{2}V_2$ from $J_2 = \frac{1}{2}\sum_e c_e\gamma_{ee}(0)$ cancels the diagonal self-term, leaving
\[
J_2 - \frac{1}{2}V_2 = \frac{1}{2}\sum_{e,e'} c_e c_{e'}\gamma_{ee'}(0) - \sum_{m=0}^{\infty}\sum_{e,e'} \gamma_{ee'}(m + 1) c_{e'}\bigl[u_{e' \to e}(m + 1) - c_e\bigr].
\]
For $e: j \to i$ and $e': j' \to i'$, edge $e$ at step $m + 1$ after $e'$ at step $0$ has probability $u_{e' \to e}(m + 1) = P(j \to i)(\mathbf{P}^m)_{j i'}$. Using $c_{e'} = w_{i'} a_{e'} v_{j'}/\lambda$, $(\mathbf{P}^m)_{j i'} = (\mathbf{M}^m)_{j i'} w_j / w_{i'}$, and $P(j \to i) = w_i a_e / (\lambda w_j)$ gives $c_{e'} u_{e' \to e}(m + 1) = \mathbf{w}^{\mathsf{T}}\mathbf{B}_e\mathbf{M}^m\mathbf{B}_{e'}\mathbf{v}$. Subtracting $c_e c_{e'} = \mathbf{w}^{\mathsf{T}}\mathbf{B}_e\mathbf{v}\mathbf{w}^{\mathsf{T}}\mathbf{B}_{e'}\mathbf{v}$ and using $\mathbf{M}^m\mathbf{Q} = \mathbf{T}_\perp^m\mathbf{Q}$ gives $c_{e'}[u_{e' \to e}(m + 1) - c_e] = \mathbf{w}^{\mathsf{T}}\mathbf{B}_e\mathbf{T}_\perp^m\mathbf{Q}\mathbf{B}_{e'}\mathbf{v}$ (with $e$ the later edge and $e'$ the earlier edge), matching $C''$ in Eq.~\eqref{eq:ring-rem-lin}.

For the remainder, class~(a) holds directly by Theorem~\ref{thm:ring-rem-lin}. For class~(b), let $\mathbf{A}_\sigma(t) = \mathbf{A}_0 + \sigma\xi(t)\mathbf{D}$ with primitive $\mathbf{A}_0$, stationary ergodic centered unit-variance $|\xi(t)| \le M_\xi$, and $\mathbf{D}$ supported on positive entries of $\mathbf{A}_0$. Choose $\sigma_1 > 0$ small enough that for $|\sigma| \le \sigma_1$ every positive entry of $\mathbf{A}_\sigma(t)$ remains at least half its baseline value, so $G_t = \mathbf{w}^{\mathsf{T}}\mathbf{A}_\sigma(t)\mathbf{y}(t) \ge \lambda/2$ pathwise. In the adapted norm $\|\cdot\|_*$ on $V_0 = \ker\mathbf{w}^{\mathsf{T}}$ from Step~2 of the proof of Theorem~\ref{thm:ring-rem-ar1}, $\|\mathbf{J}(t)\|_* \le (\lambda/G_t)[\beta + O(|\sigma|)] \le \hat{\beta} < 1$ deterministically for all paths on a smaller interval $|\sigma| \le \sigma_0$. Differentiating $\mathbf{y}(t + 1) = \mathbf{A}_\sigma(t)\mathbf{y}(t)/G_t$ up to order four therefore yields deterministic pathwise bounds $\sup_{t \ge 0, |\sigma| \le \sigma_0} |(\ln G_t)^{(j)}| \le K_j < \infty$ ($1 \le j \le 4$). Taylor's theorem at order three gives $|\E\ln G_t(\sigma) - \ln\lambda - \frac{1}{2}\sigma^2\E\ell_2(t)| \le (K_3/6)|\sigma|^3$, and if $\xi \stackrel{\mathrm{d}}{=} -\xi$ as a process then $\E\ell_3(t) = 0$ and Taylor's theorem at order four gives remainder $(K_4/24)\sigma^4$; averaging over $t$ yields Eq.~\eqref{eq:bounds-lineage}. Finally, in the non-primitive three-stage no-stay ring with one independent unit-variance rewarded transfer, at $T = 3m$ the three initial cohorts collect disjoint sums $U_1, U_2, U_3$ of $m$ independent rewards, so $\E_{\mathrm{env}}\Var_{P,\pi}(Y_{3m} \mid \text{env}) = \E[\frac{1}{3}\sum_{j=1}^3 U_j^2 - (\frac{1}{3}\sum_{j=1}^3 U_j)^2] = \frac{2}{3}m$ and $V_2 = 2/9$, whereas variance within a single cohort is identically zero.
\end{proof}
